%!TEX root = ../hutbthesis_main.tex
% 设置中文摘要
\keywordscn{MuJoCo，可变形物体，触觉反馈，实时仿真，视触融合}
\begin{abstractzh}

虚拟现实和远程交互中，可变形物体操作常常缺少力触觉反馈，实时性不足、触感不真实等状况，本文创建出一个基于MuJoCo的可变形物体实时触摸建模和仿真系统。该系统把MuJoCo物理引擎作为核心部分，加上WEART触觉手套和VR设备，形成了视觉、物理仿真、触觉反馈相结合的一体化循环系统，实现了对软体、布料、海绵等常见可变形物体的高精度建模，实时计算接触力，自动调节力反馈和动态纹理触觉显示效果，本文使用改良阻尼最小二乘法改进虚拟手逆运动学求解过程，采用多线程并行框架保证视触同步、低时延运行，通过对比实验从任务效率、操作精准度、材质识别率、用户主观感受等几个方面对系统进行验证。实验结果表明，视触融合模式相比单纯的视觉模式任务完成时间提高了30\%以上，系统端到端延迟小于15ms，可以给可变形物体交互提供稳定的、真实的、流畅的触觉仿真环境，该系统可以为虚拟手术培训、柔性物体操作、远程触觉交互等领域的标准化实验平台和技术参考提供支持，具有一定的理论价值和工程应用意义。

\end{abstractzh}
